\section*{Abstract}
The faster a drone flies through clutter, the less time it has between seeing an obstacle and reacting to it, so high-speed flight requires low perception and control latency. This latency is set by two parts of the loop: the sensor, which captures the scene only as fast as its frame rate allows, and the onboard computer, which must process each measurement in time. Small drones, moreover, can carry only limited compute and energy: a more powerful computer is heavier and draws more power, which shortens the flight and reduces agility. Neuromorphic computing offers a promising way around these limits. Inspired by biological nervous systems, it represents information as sparse spikes and computes only when they arrive, which makes it more energy-efficient than conventional deep learning. To take full advantage of it, neuromorphic computation must be coupled with event-based sensing, so that spiking neural networks (SNNs) operate directly on the sparse data the sensor produces.

Event cameras are neuromorphic vision sensors: each pixel independently and asynchronously reports changes in brightness. This gives them microsecond latency, a high dynamic range (about 140\,dB against 60\,dB for frame cameras), very little motion blur and low power consumption, which makes them well suited to fast-moving robots in challenging lighting. Their output is a sparse stream of events that an SNN can process as it arrives, so they pair naturally with neuromorphic processing.

This research aims towards fully end-to-end neuromorphic high-speed flight, which could enable agile maneuvers in challenging conditions with less compute and energy than conventional systems. The central challenge is training the SNN. Its spikes are not differentiable, and when it is trained with reinforcement learning from scratch to navigate through unseen clutter, it stalls early in the curriculum, while a conventional artificial neural network (ANN) with the same reward does not. Its continuous observations must also be encoded into spikes without leaving neurons silent. Our preliminary results in simulation, with depth images in place of events, address both problems. We distil a trained ANN teacher into a population-coded SNN student, and place the student's input receptive fields according to the teacher's observation distribution, so that almost no encoder neuron is silent. In the densest clutter, the student succeeds in 84.9\% of the flights, against 86.1\% for the teacher, at an estimated $2.3\times$ lower inference energy.

The rest of the research works towards closing the remaining gap to fully neuromorphic agile flight. We will replace the depth input with events and train the full event-to-command policy, deploy it on a neuromorphic processor, fly it on a real quadrotor, and compare it with a matched ANN in flight performance, latency and energy. We will also study whether the internal state of the SNN can serve as a memory when the onboard sensors observe only part of the state. The expected contribution is a method to train and deploy end-to-end neuromorphic policies for agile flight from onboard vision, together with a measured assessment of what they gain over ANNs.
